\subsection{Random Forest}
\label{sec:random_forest}

Model ensemble RF mengagregasikan prediksi sekumpulan pohon keputusan independen untuk mengklasifikasikan fitur wajah berdasarkan prinsip bootstrap aggregating (bagging) dan random subspace method. Pendekatan ini mereduksi varians prediksi dengan membangun sebanyak $B$ pohon klasifikasi dari sampel acak data latih. Setiap percabangan node membatasi pemilihan fitur pada subset acak berukuran max\_features $\in \{\text{'sqrt'}, \text{'log2'}\}$ untuk menekan korelasi antarpohon. Kriteria pemisahan cabang mengoptimalkan indeks Gini impurity hingga memenuhi ambang batas minimum pemisahan node min\_samples\_split, batas kedalaman max\_depth, atau batas minimum daun min\_samples\_leaf. Mekanisme majority voting menentukan keputusan akhir sampel citra melintasi seluruh pohon keputusan. Prosedur penalaan hyperparameter mengevaluasi kombinasi penskalaan fitur, reduksi dimensionalitas PCA, variasi jumlah pohon n\_estimators, serta batasan pertumbuhan pohon, menghasilkan 288 konfigurasi grid dengan 1.440 total fits melalui 5-Fold Stratified Cross-Validation sebagaimana dirangkum pada Table~\ref{tab:III}.

\begin{table}[htbp]
\caption{Hyperparameter Search Space for Random Forest Classifier.}
\label{tab:III}
\centering
\begin{tabular}{llc}
\toprule
Component / Hyperparameter & Evaluated Values & Count \\
\midrule
Feature Scaler & None, MinMaxScaler & 2 \\
Dimensionality Reduction (PCA) & None, 0.50, 0.75 & 3 \\
Number of Estimators (n\_estimators) & 100, 200 & 2 \\
Maximum Tree Depth (max\_depth) & None, 20, 30 & 3 \\
Feature Subspace (max\_features) & 'sqrt', 'log2' & 2 \\
Minimum Samples Split (min\_samples\_split) & 2, 5 & 2 \\
Minimum Samples Leaf (min\_samples\_leaf) & 1, 2 & 2 \\
\midrule
\textbf{Total Grid Combinations} & \textbf{2 × 3 × 2 × 3 × 2 × 2 × 2} & \textbf{288 (1,440 fits)} \\
\bottomrule
\end{tabular}
\end{table}
